\documentclass[a4paper,11pt,twoside]{report}
\usepackage[utf8]{inputenc}
\usepackage[T1]{fontenc}
\usepackage{lmodern}
\usepackage[margin=2.5cm]{geometry}
\usepackage{graphicx}
\usepackage{hyperref}
\usepackage{xcolor}
\usepackage{listings}
\usepackage{booktabs}
\usepackage{fancyhdr}
\usepackage{titlesec}
\usepackage{enumitem}
\usepackage{tcolorbox}

\definecolor{privaiblue}{HTML}{1a2332}
\definecolor{privaigold}{HTML}{c8a85c}
\definecolor{codebg}{HTML}{f5f5f0}
\definecolor{linkblue}{HTML}{2563eb}

\hypersetup{colorlinks=true,linkcolor=privaiblue,urlcolor=linkblue,citecolor=privaiblue}

\lstset{
  backgroundcolor=\color{codebg},
  basicstyle=\ttfamily\small,
  breaklines=true,
  frame=single,
  rulecolor=\color{gray!40},
  xleftmargin=1em,
  framexleftmargin=0.5em,
  columns=fullflexible
}

\tcbuselibrary{skins,breakable}
\newtcolorbox{infobox}[1][]{colback=blue!5,colframe=privaiblue,fonttitle=\bfseries,title=#1,breakable}
\newtcolorbox{warnbox}[1][]{colback=yellow!8,colframe=privaigold,fonttitle=\bfseries,title=#1,breakable}

\pagestyle{fancy}
\fancyhf{}
\fancyhead[LE,RO]{\thepage}
\fancyhead[LO]{\nouppercase{\rightmark}}
\fancyhead[RE]{Privai User Guide}
\renewcommand{\headrulewidth}{0.4pt}

\titleformat{\chapter}[display]{\normalfont\huge\bfseries\color{privaiblue}}{Chapter \thechapter}{12pt}{\Huge}

\title{
  \vspace{2cm}
  {\Huge\bfseries\color{privaiblue} Privai}\\[0.5cm]
  {\Large Privacy-First AI Desktop Workspace}\\[1cm]
  {\large User Guide \& Technical Documentation}\\[2cm]
  {\normalsize Version 0.2.0}
}
\author{Privai Project Documentation}
\date{May 2026}

\begin{document}
\maketitle
\tableofcontents

%========================================================================
\chapter{Introduction}
%========================================================================

\section{What Is Privai?}
Privai is a \textbf{privacy-first AI desktop workspace} for business automation, document management, and application building. It runs as a native macOS desktop application (Electron) with a local FastAPI backend that communicates with the OpenAI API for language model responses.

The core privacy principle: \textbf{all chat history stays on your device} in a local SQLite database. Firebase is used \emph{only} for user identity (sign-in and ID tokens). No conversation content is stored in the cloud.

\section{Key Features}
\begin{itemize}[leftmargin=1.5em]
  \item \textbf{Chat with AI} --- conversational assistant powered by OpenAI models (GPT-5.4-mini default)
  \item \textbf{Web Search} --- integrated SearXNG meta-search with automatic source citation
  \item \textbf{Privacy Guard} --- regex-based redaction of sensitive data (CNP, phone, email, API keys) before it leaves the device
  \item \textbf{Agent Mode} --- terminal-based code agent that can inspect, create, and edit files in a workspace
  \item \textbf{File Conversion} --- local image/document conversion (HEIC, PDF, etc.)
  \item \textbf{PDF Generation} --- create polished PDF documents from natural language
  \item \textbf{Learning Notebooks} --- upload study materials and quiz yourself
  \item \textbf{Google Workspace} --- read-only Gmail search and Google Calendar integration
  \item \textbf{Workspaces} --- four specialized spaces: General, Business, Coding, and Learning
  \item \textbf{Streaming Responses} --- real-time SSE token streaming
  \item \textbf{Artifact Rendering} --- agent-generated HTML rendered in a sandboxed preview
\end{itemize}

\section{Architecture Overview}
The system consists of three components:

\begin{enumerate}
  \item \textbf{FastAPI Backend} (Python, port 8080) --- orchestrator, LLM client, search, privacy guard, SQLite, terminal tools
  \item \textbf{Web App} (Next.js 16, port 3000/3100) --- polished chat UI with Firebase Auth
  \item \textbf{Agent Service} (Node.js) --- Firestore heartbeat writer for device online/offline status
\end{enumerate}

\begin{infobox}[Privacy Posture]
Chat messages, sources, and redactions live \textbf{only} in local SQLite (\texttt{data/chat.db}). Firebase stores only authentication credentials and one heartbeat document per device. The OpenAI API sees prompt content for generation, but nothing is persisted on their side (\texttt{store: false}).
\end{infobox}

%========================================================================
\chapter{Installation \& Setup}
%========================================================================

\section{Prerequisites}
\begin{itemize}
  \item macOS (Apple Silicon supported), Python 3.11+, Node.js 20+, Git
  \item An OpenAI API key
\end{itemize}

\section{Quick Install}
\begin{lstlisting}[language=bash]
# Clone the repository, then:
bash scripts/install.sh
\end{lstlisting}
This creates \texttt{.venv/}, installs FastAPI dependencies, clones SearXNG, and installs its requirements. It also creates \texttt{data/} directories and copies \texttt{.env.example} to \texttt{.env}.

\section{Configuration}
Edit the \texttt{.env} file in the project root:
\begin{lstlisting}
OPENAI_API_KEY=your-openai-api-key
OPENAI_MODEL=gpt-5.4-mini
OPENAI_VISION_MODEL=gpt-5.4-mini
\end{lstlisting}

\subsection{All Environment Variables}
\begin{center}\small
\begin{tabular}{@{}llp{5.5cm}@{}}
\toprule
\textbf{Variable} & \textbf{Default} & \textbf{Purpose} \\
\midrule
\texttt{OPENAI\_API\_KEY} & --- & OpenAI API key \\
\texttt{OPENAI\_MODEL} & \texttt{gpt-5.4-mini} & Default text model \\
\texttt{OPENAI\_VISION\_MODEL} & \texttt{gpt-5.4-mini} & Model for image attachments \\
\texttt{OPENAI\_REASONING\_EFFORT} & \texttt{low} & Reasoning effort for reasoning models \\
\texttt{SEARXNG\_URL} & \texttt{http://127.0.0.1:8888} & SearXNG base URL \\
\texttt{API\_HOST} & \texttt{0.0.0.0} & Bind address \\
\texttt{API\_PORT} & \texttt{8080} & Bind port \\
\texttt{SEARCH\_TOP\_K} & \texttt{50} & Search results passed to LLM \\
\texttt{SEARCH\_FALLBACK\_ENABLED} & \texttt{true} & DuckDuckGo fallback when SearXNG is offline \\
\texttt{TERMINAL\_ENABLED} & \texttt{true} & Enable agent terminal tools \\
\texttt{TERMINAL\_TIMEOUT\_S} & \texttt{60} & Command timeout in seconds \\
\texttt{TERMINAL\_ALLOW\_DANGEROUS} & \texttt{false} & Bypass destructive command denylist \\
\texttt{AGENT\_MAX\_TOOL\_STEPS} & \texttt{20} & Max tool loop steps per agent reply \\
\bottomrule
\end{tabular}
\end{center}

\section{Running the Full Stack}
Three processes are needed:

\begin{lstlisting}[language=bash]
# 1) Backend + SearXNG (prints PAIRING CODE on first boot)
bash scripts/run_all.sh

# 2) Agent heartbeat service
cd agent && cp .env.example .env
npm install && npm run dev

# 3) Web app
cd web && cp .env.local.example .env.local
npm install && npm run dev    # http://localhost:3000
\end{lstlisting}

\section{Desktop Application}
The Electron shell is in \texttt{web/electron/}. It auto-starts the FastAPI backend and the Next.js UI.

\begin{lstlisting}[language=bash]
# Run from source:
cd web && npm run desktop:dev

# Build distributable:
cd web && npm run dist
\end{lstlisting}

Installed desktop builds read API keys from:\\
\texttt{\textasciitilde/Library/Application Support/Privai/.env}

%========================================================================
\chapter{Device Pairing}
%========================================================================

\section{How Pairing Works}
On first boot (when no owner is set), the backend prints a 6-digit pairing code to the console:

\begin{lstlisting}
+----------------------------------------------+
|  PAIRING CODE: 482919                        |
|  Enter this in the web app /settings page.   |
+----------------------------------------------+
\end{lstlisting}

\section{Pairing Steps}
\begin{enumerate}
  \item Start the backend with \texttt{bash scripts/run\_all.sh}
  \item Open \texttt{http://localhost:3000} and sign in (Email, Google, or GitHub)
  \item Go to \textbf{Settings}
  \item Enter the 6-digit code and click \emph{Pair device}
  \item The device is now bound to your Firebase UID
\end{enumerate}

\begin{infobox}[Re-pairing]
To pair a different account: delete \texttt{data/chat.db}, restart the backend, and repeat the pairing flow.
\end{infobox}

%========================================================================
\chapter{Using Privai}
%========================================================================

\section{Workspaces (Spaces)}
Privai organizes conversations into four spaces:

\begin{description}[leftmargin=2em]
  \item[General] Default space for everyday questions and tasks.
  \item[Business] Business automation, email search, calendar management, client workflows.
  \item[Coding] Software development with agent terminal tools, file editing, and app building.
  \item[Learning] Study notebooks with material upload, quizzing, and key-term extraction.
\end{description}

\section{Chat Mode}
The default mode. Type a question and receive an AI-generated answer. Features:
\begin{itemize}
  \item \textbf{Automatic web search}: triggered by keywords like ``search'', ``latest'', ``news'', ``price'', URLs, etc.
  \item \textbf{Manual web toggle}: force search on/off per message.
  \item \textbf{Source citations}: search results are numbered \texttt{[1]}, \texttt{[2]}, etc.\ in the response.
  \item \textbf{Streaming}: responses arrive token-by-token via SSE.
  \item \textbf{Image attachments}: upload images for vision model analysis.
  \item \textbf{File attachments}: text files are embedded as context in the prompt.
\end{itemize}

\section{Agent Mode}
Press the \textbf{Agent} toggle to enter agent mode. The AI can:
\begin{itemize}
  \item Run shell commands in the selected workspace
  \item Create and edit project files
  \item Install dependencies and run builds
  \item Verify results with lint, tests, or build checks
\end{itemize}

\begin{warnbox}[Command Approval]
By default, each terminal command requires your approval before execution. Read-only commands (ls, cat, grep, etc.) can be auto-approved. Dangerous commands (sudo, rm -rf /, etc.) are blocked unless \texttt{TERMINAL\_ALLOW\_DANGEROUS=true}.
\end{warnbox}

\noindent Example prompts for agent mode:
\begin{lstlisting}
Inspect this repo and tell me how to run the tests.
Create a new Vite React app called todo-demo.
Search the backend for streaming code and explain the flow.
\end{lstlisting}

Use \textbf{File > Open Workspace...} in the desktop app to choose where the agent operates.

\section{Convert Mode}
Attach files and ask Privai to convert them. Supported conversions:
\begin{itemize}
  \item \textbf{Image formats}: PNG, JPEG, HEIC, WebP, TIFF, BMP, GIF, ICO (via macOS \texttt{sips})
  \item \textbf{Office to PDF}: DOCX, PPTX, XLSX, ODT, etc.\ (via LibreOffice \texttt{soffice})
  \item \textbf{AI-generated PDFs}: ask ``create a PDF about X'' and the LLM generates the content
\end{itemize}

\section{Learning Notebooks}
In the Learning space:
\begin{enumerate}
  \item Create a learning session
  \item Add materials: paste text or upload files (PDF, DOCX, images)
  \item The system extracts text, identifies key terms, and creates summaries
  \item Ask questions --- the AI answers from your materials first, using general knowledge only when needed
\end{enumerate}

\section{Google Workspace Integration}
In the Business space, connect your Google account for:
\begin{itemize}
  \item \textbf{Gmail search} (read-only): search and read email threads
  \item \textbf{Calendar}: find free slots and draft calendar events
  \item Calendar events are created as \emph{pending Business actions} requiring your approval
\end{itemize}

Setup: Go to Settings, click ``Connect Google Workspace'', and complete the OAuth flow.

\section{Session Management}
\begin{itemize}
  \item Sessions are auto-titled by the LLM after the first message
  \item Rename, delete, or clear sessions from the sidebar
  \item \textbf{Compact}: summarize a long conversation to save context window space
  \item \textbf{Context meter}: shows approximate token usage relative to the model's context window
\end{itemize}

%========================================================================
\chapter{Privacy Guard}
%========================================================================

The Privacy Guard is a regex-based filter that scans user input \textbf{before} it leaves the device (e.g., before being sent to SearXNG for web search). It redacts:

\begin{center}\small
\begin{tabular}{@{}ll@{}}
\toprule
\textbf{Pattern} & \textbf{Description} \\
\midrule
CNP & Romanian personal identification number (13 digits) \\
EMAIL & Email addresses \\
PHONE & Romanian mobile numbers (\texttt{+40 7xx...}) \\
CARD & Credit card-like digit sequences (13--19 digits) \\
APIKEY & API keys, tokens, bearer tokens \\
PASSWORD & Lines matching \texttt{password: ...} patterns \\
\bottomrule
\end{tabular}
\end{center}

Redacted content appears as \texttt{[REDACTED:LABEL]} in the sanitized text. The UI shows a yellow tag when redactions occur, demonstrating to the user that sensitive data was stripped.

%========================================================================
\chapter{API Reference}
%========================================================================

All endpoints except \texttt{/health}, \texttt{/pair}, and \texttt{/pair/status} require a Firebase ID token as \texttt{Authorization: Bearer <token>}.

\begin{center}\small
\begin{tabular}{@{}llll@{}}
\toprule
\textbf{Method} & \textbf{Path} & \textbf{Auth} & \textbf{Description} \\
\midrule
GET & \texttt{/health} & --- & System health and status \\
POST & \texttt{/pair} & token & Claim device with pairing code \\
GET & \texttt{/pair/status} & --- & Check pairing status \\
\midrule
GET & \texttt{/sessions} & owner & List chat sessions \\
POST & \texttt{/sessions} & owner & Create a new session \\
GET & \texttt{/sessions/\{id\}} & owner & Get one session \\
PATCH & \texttt{/sessions/\{id\}} & owner & Rename session \\
DELETE & \texttt{/sessions/\{id\}} & owner & Delete session \\
GET & \texttt{/sessions/\{id\}/messages} & owner & List messages \\
\midrule
POST & \texttt{/chat} & owner & Non-streaming chat \\
POST & \texttt{/chat/stream} & owner & SSE streaming chat \\
POST & \texttt{/search} & owner & Raw SearXNG search \\
\midrule
GET & \texttt{/admin/llm/status} & owner & Model status \\
GET & \texttt{/admin/llm/models} & owner & Available models \\
POST & \texttt{/admin/llm/model} & owner & Change active model \\
\midrule
POST & \texttt{/attachments} & owner & Upload file \\
GET & \texttt{/attachments/\{id\}/raw} & owner & Download file \\
DELETE & \texttt{/attachments/\{id\}} & owner & Delete attachment \\
\midrule
GET & \texttt{/workspace/tree} & owner & Browse workspace \\
GET & \texttt{/workspace/file} & owner & Read a file \\
POST & \texttt{/workspace/file} & owner & Save a file \\
POST & \texttt{/workspace/terminal} & owner & Run terminal command \\
\bottomrule
\end{tabular}
\end{center}

%========================================================================
\chapter{Project Structure}
%========================================================================

\begin{lstlisting}
backend/        FastAPI app, orchestrator, privacy guard,
                LLM client, SQLite, terminal tools
frontend/       Static HTML/CSS/JS chat UI fallback
web/            Next.js 16 + Tailwind + Firebase Auth
  electron/     Electron shell for desktop app
  src/app/      Next.js page routes
  src/components/ React UI components
  src/lib/      Utilities, types, Firebase config
agent/          Node.js heartbeat service
searxng/        SearXNG upstream clone (gitignored)
scripts/        install.sh, run_all.sh, run_api.sh, etc.
data/           chat.db, attachments, cache, logs
docs/           Documentation and assets
firestore.*     Firestore rules and indexes
firebase.json   Firebase Hosting config
\end{lstlisting}

\subsection{Key Backend Modules}
\begin{description}[leftmargin=2em,style=nextline]
  \item[\texttt{main.py}] FastAPI routes (1683 lines): chat, sessions, attachments, workspace, learning, business actions.
  \item[\texttt{orchestrator.py}] Decides search need, builds LLM prompts, manages agent tool loops.
  \item[\texttt{llm\_client.py}] Async OpenAI Responses API client with streaming and tool calling.
  \item[\texttt{privacy\_guard.py}] Regex scanner that redacts sensitive patterns.
  \item[\texttt{database.py}] SQLite layer: sessions, messages, attachments, learning materials, business actions, KV store.
  \item[\texttt{auth.py}] Firebase ID token verification using Google's public x509 certificates.
  \item[\texttt{searxng\_client.py}] Async SearXNG JSON search with DuckDuckGo HTML fallback.
  \item[\texttt{terminal.py}] Sandboxed shell command execution with timeout, output capture, and safety checks.
  \item[\texttt{converter.py}] Local file conversion using \texttt{sips} and LibreOffice.
  \item[\texttt{pdf\_builder.py}] Pure-Python PDF renderer from Markdown with image embedding.
  \item[\texttt{google\_workspace.py}] OAuth2, Gmail read-only, Google Calendar integration.
  \item[\texttt{artifacts.py}] Extracts \texttt{<artifact>} HTML blocks from agent responses.
  \item[\texttt{approvals.py}] In-memory approval gate for agent terminal commands.
\end{description}

%========================================================================
\chapter{Troubleshooting}
%========================================================================

\begin{description}[leftmargin=1em]
  \item[``OPENAI\_API\_KEY is not set''] Add your key to \texttt{.env} or \texttt{\textasciitilde/Library/Application Support/Privai/.env}.
  \item[SearXNG not starting] Ensure \texttt{searxng/searx/settings\_local.yml} exists. Run \texttt{bash scripts/install.sh} again.
  \item[Pairing code not appearing] The device is already paired. Check \texttt{GET /pair/status}.
  \item[Agent commands blocked] Dangerous commands are blocked by default. Set \texttt{TERMINAL\_ALLOW\_DANGEROUS=true} to override.
  \item[Google Workspace errors] Ensure \texttt{GOOGLE\_CLIENT\_ID} and \texttt{GOOGLE\_CLIENT\_SECRET} are set. Complete the OAuth flow in Settings.
  \item[Context too long] Use the Compact button to summarize the conversation, or start a new session.
  \item[File conversion fails] Image conversion requires macOS \texttt{sips}. Office-to-PDF requires LibreOffice (\texttt{brew install --cask libreoffice}).
\end{description}

%========================================================================
\chapter*{Quick Reference Card}
\addcontentsline{toc}{chapter}{Quick Reference Card}
%========================================================================

\begin{center}
\begin{tabular}{@{}p{4cm}p{9cm}@{}}
\toprule
\textbf{Action} & \textbf{How} \\
\midrule
Start everything & \texttt{bash scripts/run\_all.sh} \\
Open web UI & \texttt{http://localhost:3000} \\
Desktop app (dev) & \texttt{cd web \&\& npm run desktop:dev} \\
Build installer & \texttt{cd web \&\& npm run dist} \\
Change model & Settings page or \texttt{POST /admin/llm/model} \\
Force web search & Toggle the search icon in the composer \\
Use agent mode & Toggle the Agent button in the composer \\
Convert files & Attach files + type ``convert to PDF'' \\
Open workspace & File > Open Workspace (desktop app) \\
Pair new device & Delete \texttt{data/chat.db}, restart backend \\
\bottomrule
\end{tabular}
\end{center}

\end{document}
